\section{Introduction}

A popular way to frame the renewable-energy question is ``how much of the Earth
would we need to cover in solar panels?'' The framing is seductive because the
raw numbers are staggering: the Sun deposits roughly \InterceptedTW~TW on
Earth's cross-section, against a human primary-energy appetite near 18~TW. By
that arithmetic the problem is trivially solved and the only question is cost.

This paper argues that the arithmetic hides the real constraint. The
\InterceptedTW~TW is not idle: after the atmosphere reflects and absorbs its
share, the sunlight that reaches the ground is what powers photosynthesis,
evaporation, and every ecosystem on the planet. Harvesting solar energy at
civilisational scale is therefore not a question of \emph{area} but of
\emph{allocation} --- how much light can we divert to electricity before we
starve the biosphere that light supports?

We approach this with a deliberately self-contained, first-principles model
(no external datasets; every number is reproducible from published physical
constants) organised around four questions:

\begin{enumerate}
\item \textbf{The ceiling.} How much sunlight actually reaches the surface,
  and where does the rest go? We track the cascade from the solar constant to
  surface irradiance, breaking out the ozone layer's contribution, and map the
  resource by latitude band (Section~\ref{sec:resource}).
\item \textbf{Feasibility.} How much energy is deployable on land without
  displacing forests, croplands, or protected habitat, and how does that
  compare to demand? (Section~\ref{sec:feasibility}).
\item \textbf{The orbital alternative.} Is a collector in space --- a ``layer''
  above the atmosphere rather than a Dyson sphere --- fundamentally better, and
  why? (Section~\ref{sec:feasibility}).
\item \textbf{Sharing light with life.} If an ecosystem needs a certain amount
  of light but is flexible about \emph{when} it receives it, how much more can
  we harvest by scheduling around it? (Section~\ref{sec:ecosystem}).
\end{enumerate}

The fourth question is our main contribution. We formalise ecosystem-aware
harvesting as a constrained optimisation, prove its optimum is a simple
water-filling rule, and quantify the energy unlocked purely by respecting
\emph{when} --- not just how much --- light a biome needs.
